\documentclass[11pt,a4paper]{article}

\usepackage{fontspec}
\usepackage[ngerman]{babel}
\usepackage[a4paper,top=27mm,bottom=27mm,left=30mm,right=30mm]{geometry}
\usepackage{microtype}
\usepackage{graphicx}
\usepackage{xcolor}
\usepackage{titlesec}
\usepackage{enumitem}
\usepackage{parskip}
\usepackage{fancyhdr}
\usepackage{tcolorbox}

\definecolor{tiefblau}{HTML}{16324F}
\definecolor{stahl}{HTML}{3B6EA5}
\definecolor{sand}{HTML}{F6F2EA}
\definecolor{grau}{HTML}{5A6673}
\definecolor{rostrot}{HTML}{9C4221}

\setmainfont{Palatino}[
  Numbers = OldStyle,
  Ligatures = TeX
]
\newfontfamily\displayfont{Helvetica Neue}[Ligatures=TeX]

\titleformat{\section}{\displayfont\large\bfseries\color{tiefblau}}{}{0pt}{}
\titlespacing*{\section}{0pt}{18pt}{6pt}

\pagestyle{fancy}
\fancyhf{}
\renewcommand{\headrulewidth}{0pt}
\fancyfoot[C]{\displayfont\footnotesize\color{grau}\thepage}

\newtcolorbox{befund}{
  colback=sand, colframe=rostrot, boxrule=0.8pt, arc=2pt,
  left=10pt, right=10pt, top=8pt, bottom=8pt
}

\begin{document}

\begin{center}
  {\displayfont\Huge\bfseries\color{tiefblau} brainlehr\par}
  \vspace{4pt}
  {\displayfont\large\color{stahl} Was es sein möchte, und was es schon ist\par}
  \vspace{6pt}
  {\displayfont\footnotesize\color{grau} Gemessen am 2026-08-16 · Zweig \texttt{brainlehr/b4-ausweis}\par}
\end{center}

\vspace{10pt}
\hrule height 0.6pt
\vspace{14pt}

\section{Was brainlehr sein möchte}

brainlehr versteht sich nicht in erster Linie als Wissensspeicher, sondern als
\textbf{Aufsicht über die eigene Arbeitsweise}: eine Instanz, die prüft, ob das,
was man sich vornimmt, tatsächlich wirkt --- nicht nur gebaut und gemeldet wird.
Der Anlass dafür war ein Befund von brainlehr über brainlehr: zwölf Fälle
„gebaut, laufend, meldend, wirkungslos“ an einem einzigen Tag.

Architektonisch trägt brainlehr die untere von zwei Schichten im Verbund:
\emph{was gilt, und ob es belegt ist} --- Ausweis, Freigabe, Norm mit Rang und
Geltung, Herkunft, Widerruf. Das Unterscheidungskriterium zur oberen Schicht ist
bewusst schmal: \textbf{brainlehr kann nein sagen.} Eine Freigabe wird verweigert,
eine Norm greift, ein Trigger blockiert, ein Melder schlägt an. Die fachlichen
Instanzen können das nicht --- sie stellen bereit.

Vier Fragen soll jede Instanz beantworten: Wer fragt hier? Worüber wird Wissen
geführt? Was ist ein Treffer wert? Und was darf nach außen?

\vspace{6pt}
\begin{center}
  \includegraphics[width=\textwidth]{graph.png}
\end{center}
\vspace{2pt}

\section{Was es schon ist --- gemessen}

\textbf{Bestand} (\texttt{brainlehr.db}, 32 Tabellen): 5034 Knoten, 963 Lehren,
9936 Kanten, 16\,144 protokollierte Zugriffe, 55 Datenbank-Trigger.

Diese Zahlen liegen deutlich über den in der eigenen Projektbeschreibung
genannten (2166 Knoten, 833 Lehren). Die Beschreibung veraltet also schneller,
als sie fortgeschrieben wird --- ein Beispiel für genau das Muster, das das
System an anderer Stelle sucht.

\textbf{Gebaut:} 33 Melder-Skripte, 95 Dateien unter \texttt{kern/},
224 Testdateien, 24 Architekturentscheidungen, 99 Dokumente, davon 48 Pläne.

\begin{befund}
\textbf{Der zentrale Befund --- und er hat sich beim Aufschreiben selbst
korrigiert.}

Der Melder meldete zunächst \textbf{29 von 33} Mechanismen ohne Auslöser. Beim
Gegenlesen fiel auf, dass darunter sechs standen, die im \texttt{pre-push}
hängen und an diesem Tag mehrfach einen Push \emph{tatsächlich} gestoppt haben.
Ursache: Der Melder kannte Einstellungen, geplante Läufe und die Aufruferkette
--- aber keine Git-Hooks. Ein Eintrag in den Einstellungen \emph{löst aus}, ein
Git-Hook \emph{blockiert}; wer den zweiten nicht kennt, meldet ausgerechnet die
wirksamsten Mechanismen als wirkungslos.

Nach der Korrektur: \textbf{21 von 33 ohne Auslöser, 12 verdrahtet.} Das bleibt
der harte Teil des Befunds --- aber die Zahl stimmt jetzt.
\end{befund}

\textbf{Was nachweislich wirkt:} Der \texttt{pre-push}-Wächter prüft
Messauswertungen gegen die Messregeln und vergleicht die Landkarten gegen eine
Neuerzeugung. Er blockiert bei einem Verstoß tatsächlich --- an einem
Rollout wurde er zweimal falsch positiv grün gemessen; erst der dritte, korrekte
Aufbau zeigte, dass er greift.

\textbf{Offene Stellen:} Die Relevanzschwelle ist für einen Funktionsnamen und
einen Rechtssatz identisch, obwohl der Unterschied als Lücke benannt ist. Der
Bedeutungskanal der Suche steuert nach eigener Messung nur 4 von 585 Endplätzen
bei; der Stichwortkanal dominiert. Eine gebaute Fusionsfunktion lag seit dem
2026-08-12 unverdrahtet im Code.

\vspace{10pt}
\hrule height 0.4pt
\vspace{10pt}

{\displayfont\normalsize\color{tiefblau}\textbf{Fazit in einem Satz}}

Der Anspruch --- Aufsicht statt Speicher --- trägt: Die Selbstprüfung hat beim
Schreiben dieses Papiers einen Fehler in sich selbst gefunden und korrigiert. An
21 von 33 eigenen Mechanismen ist er trotzdem noch Absicht ohne Auslöser.

\vspace{16pt}
{\displayfont\footnotesize\color{grau}
Jede Zahl in diesem Papier wurde für dieses Papier gemessen, nicht übernommen.
Die Zählweise steht im Erhebungsprotokoll.\par}

\end{document}
